\documentclass[11pt,oneside,notitlepage]{ctexrep}

\usepackage[a4paper,margin=1in]{geometry}
\usepackage{amsmath,amssymb,amsthm,mathtools,booktabs}
\usepackage{enumitem}
\usepackage{tikz}
\usepackage[hidelinks]{hyperref}

\ctexset{chapter/format=\centering\Large\bfseries}

\theoremstyle{definition}
\newtheorem{definition}{定义}[chapter]
\newtheorem{proposition}[definition]{命题}
\newtheorem{lemma}[definition]{引理}
\newtheorem{corollary}[definition]{推论}
\newtheorem{theorem}[definition]{定理}
\theoremstyle{remark}
\newtheorem{remark}[definition]{注}

\newcommand{\cJ}{\mathcal{J}}
\newcommand{\cI}{\mathcal{I}}
\newcommand{\cK}{\mathcal{K}}
\newcommand{\Ustar}{U^{*}}

% 甘特图：\gbar{行}{开始}{结束}{标签}{颜色}
\newcommand{\gbar}[5]{%
  \fill[#5] (#2,-#1-0.35) rectangle (#3,-#1+0.35);
  \draw (#2,-#1-0.35) rectangle (#3,-#1+0.35);
  \node[font=\small] at ({(#2+#3)/2},-#1) {#4};}
\newcommand{\gaxis}[2]{% 最后一行的编号, 最大时刻；时间轴画在最后一行下方
  \draw[->] (0,-#1-0.55) -- (#2+0.4,-#1-0.55);
  \foreach \t in {0,...,#2} {\draw (\t,-#1-0.55) -- (\t,-#1-0.65); \node[font=\scriptsize,below] at (\t,-#1-0.6) {\t};}}

\title{毕业论文第2章\quad 问题描述与数学模型（草稿）}
\author{FJSP-AGV-RL Research Platform}
\date{2026-11-01}

\begin{document}
\maketitle

\paragraph{草稿说明}
本文是论文第2章的草稿，对应仓库当前实现（\texttt{src/fjsp\_lab/maintenance/}）。
模型假设来自用户确认的毕设首版设定，见 \texttt{docs/models/graduation\_model.md}。
各结论的状态如下：
\begin{itemize}[nosep]
  \item 命题~\ref{prop:fixed-joint}、引理~\ref{lem:left-shift}、推论~\ref{cor:horizon}、
        命题~\ref{prop:env} 在本文给出完整证明；
  \item 定理~\ref{thm:np} 的归约由本文推导，所引用的经典结果见参考文献，\textbf{需作者逐步核对}；
  \item 注~\ref{rem:decoder} 所述的解码器构造空间能否覆盖最优解\textbf{没有证明}，
        只在小规模算例上穷举验证过。
\end{itemize}

\setcounter{chapter}{1}
\chapter{问题描述与数学模型}

\section{问题描述}

考虑由 $m$ 台同质平行机组成的车间，需要加工 $n$ 个工件。
每个工件只有一道工序，可以在任意一台机器上加工，加工时间与机器无关；
所有工件在 0 时刻到达，各有一个交货期。
每台机器有若干个预先给定、按时间排列且互不重叠的预防维护时间窗；
每个时间窗内必须完整地做一次维护，维护期间该机器不能加工。
即使某台机器没有分到任何工件，它的维护也必须执行。
目标是最小化总延期。

调度需要同时决定三件事：工件分配到哪台机器（任务分配）、每台机器上工件的先后（加工顺序）、
每次维护的开始时刻（维护时机）。本文比较两种维护方式：
\begin{itemize}[nosep]
  \item \textbf{固定维护}：每次维护的开始时刻预先给定，调度只决定任务分配和加工顺序；
  \item \textbf{联合维护}（协同优化）：维护开始时刻也作为决策变量，在时间窗内自由选择。
\end{itemize}

每台机器只有一个时间窗时，称为\textbf{单窗口模型}，它对应 Lee 与 Chen~\cite{lee2000}
研究的“每台机器在给定期限内恰好维护一次、不同机器可同时维护”的情形；
本文把期限 $[0,T]$ 推广为一般时间窗 $[L,U)$，目标由加权总完工时间改为总延期，
并增加了固定与联合两种维护方式的对比。
允许每台机器有多个时间窗、且各机器窗口数可以不同时，称为\textbf{周期性多窗口模型}。
单窗口模型是多窗口模型的特例，下文统一按多窗口模型叙述。

模型采用以下假设：
\begin{enumerate}[label=(A\arabic*),nosep]
  \item 机器同质，工件 $j$ 在任一机器上的加工时间都是 $p_j$；
  \item 每个工件一道工序，加工不可中断，工件之间没有先后约束，全部在 0 时刻到达；
  \item 机器 $i$ 的每个时间窗内必须恰好执行一次维护，维护时长给定，维护不能省略；
  \item 整个维护区间必须落在时间窗内，而不是只约束维护的开始时刻；
  \item 维护与该机器上的加工互斥；不同机器的维护可以同时进行，不考虑维修人员数量限制；
  \item 不考虑运输、换型准备、随机故障和抢占；所有时间为非负整数。
\end{enumerate}

\section{符号}

\begin{table}[h]
\centering
\caption{集合、参数与决策变量}
\label{tab:notation}
\begin{tabular}{ll}
\toprule
符号 & 含义 \\
\midrule
$\cJ=\{1,\dots,n\}$ & 工件集合，$n\ge 1$ \\
$\cI=\{1,\dots,m\}$ & 机器集合，$m\ge 1$ \\
$\cK_i=\{1,\dots,K_i\}$ & 机器 $i$ 的维护时间窗集合，$K_i\ge 1$ \\
$W=\sum_{i}K_i$ & 时间窗（即维护）总数 \\
\midrule
$p_j\in\mathbb{Z}_{\ge 1}$ & 工件 $j$ 的加工时间 \\
$d_j\in\mathbb{Z}_{\ge 0}$ & 工件 $j$ 的交货期 \\
$q_{ik}\in\mathbb{Z}_{\ge 1}$ & 机器 $i$ 第 $k$ 次维护的时长 \\
$[L_{ik},U_{ik})$ & 机器 $i$ 第 $k$ 个维护时间窗（左闭右开） \\
$f_{ik}$ & 固定维护方式下预先给定的维护开始时刻 \\
$P=\sum_j p_j$ & 总加工时间 \\
$\Ustar=\max_{i,k}U_{ik}$ & 最晚的时间窗右端点 \\
$H=\Ustar+P$ & 时域上界（见推论~\ref{cor:horizon}） \\
\midrule
$x_{ji}\in\{0,1\}$ & 工件 $j$ 是否分配到机器 $i$ \\
$S_j,\;C_j$ & 工件 $j$ 的开工与完工时刻，$C_j=S_j+p_j$ \\
$M_{ik}$ & 机器 $i$ 第 $k$ 次维护的开始时刻 \\
$T_j$ & 工件 $j$ 的延期，$T_j=\max(0,\,C_j-d_j)$ \\
\bottomrule
\end{tabular}
\end{table}

本章符号见表~\ref{tab:notation}。输入数据满足
\begin{gather}
0\le L_{ik},\qquad q_{ik}\le U_{ik}-L_{ik}, \label{eq:data-window}\\
U_{ik}\le L_{i,k+1}\quad (k<K_i), \label{eq:data-order}\\
L_{ik}\le f_{ik}\le U_{ik}-q_{ik}\quad\text{（仅固定维护方式）}. \label{eq:data-fixed}
\end{gather}
式~\eqref{eq:data-window} 保证每次维护都能放进其时间窗，
式~\eqref{eq:data-order} 表示同一机器的时间窗按时间排列且互不重叠，
式~\eqref{eq:data-fixed} 保证预先给定的维护时刻本身合法。
加工与维护都占用左闭右开区间：工件 $j$ 占用 $[S_j,C_j)$，维护占用 $[M_{ik},M_{ik}+q_{ik})$。
两个区间 $[a,b)$ 与 $[c,e)$ 不重叠当且仅当 $b\le c$ 或 $e\le a$。

\section{数学模型}

联合维护方式下的模型记为 (JM)：
\begin{align}
\min\quad & \sum_{j\in\cJ} T_j \label{eq:obj}\\
\text{s.t.}\quad
& \sum_{i\in\cI} x_{ji}=1, && j\in\cJ, \label{eq:assign}\\
& C_j=S_j+p_j,\quad S_j\ge 0, && j\in\cJ, \label{eq:proc}\\
& T_j=\max(0,\,C_j-d_j), && j\in\cJ, \label{eq:tard}\\
& L_{ik}\le M_{ik}\le U_{ik}-q_{ik}, && i\in\cI,\ k\in\cK_i, \label{eq:window}\\
& x_{ji}=x_{li}=1\ \Rightarrow\ C_j\le S_l\ \lor\ C_l\le S_j, && j\ne l,\ i\in\cI, \label{eq:jj}\\
& x_{ji}=1\ \Rightarrow\ C_j\le M_{ik}\ \lor\ M_{ik}+q_{ik}\le S_j, && j\in\cJ,\ i\in\cI,\ k\in\cK_i, \label{eq:jm}\\
& x_{ji}\in\{0,1\},\quad S_j,\,M_{ik}\in\mathbb{Z}. && \label{eq:domain}
\end{align}
式~\eqref{eq:assign} 要求每个工件恰好分配到一台机器；
式~\eqref{eq:proc} 表示加工不可中断；
式~\eqref{eq:tard} 是延期的精确定义，提前完工不产生负延期，也不抵消其他工件的延期；
式~\eqref{eq:window} 要求整个维护区间落在时间窗内，对每个时间窗都必须成立，
因此维护不能省略，空闲机器也要维护；
式~\eqref{eq:jj} 和式~\eqref{eq:jm} 分别是同一机器上加工与加工、加工与维护的不重叠约束。

同一机器上两次维护之间不需要单独的约束：由式~\eqref{eq:window} 和式~\eqref{eq:data-order}，
$M_{ik}+q_{ik}\le U_{ik}\le L_{i,k+1}\le M_{i,k+1}$，两次维护自然不重叠，并且按窗口顺序执行。

固定维护方式下的模型记为 (FM)，它在 (JM) 的基础上增加
\begin{equation}
M_{ik}=f_{ik},\qquad i\in\cI,\ k\in\cK_i. \label{eq:fixed}
\end{equation}

\begin{remark}
目标 \eqref{eq:obj} 是\textbf{正则目标}：每个 $T_j$ 是 $C_j$ 的非减函数，因此任何工件提前完工都不会使目标变差。
引理~\ref{lem:left-shift} 和命题~\ref{prop:env} 都依赖这一点；
若改为提前也受罚、按时段计费等非正则目标，相关结论需要重新论证。
\end{remark}

\begin{remark}
析取约束 \eqref{eq:jj}、\eqref{eq:jm} 可以按常规方式引入排序 0--1 变量并以 $H$ 为大 $M$ 线性化，
得到混合整数线性规划；由推论~\ref{cor:horizon}，限制在 $[0,H]$ 内不损失最优值。
本文的精确求解采用第~\ref{sec:cpsat} 节的约束规划模型，不使用线性化形式。
\end{remark}

\section{固定维护与联合维护的关系}

\begin{proposition}\label{prop:fixed-joint}
记 $\mathcal{F}_{\mathrm{FM}}$、$\mathcal{F}_{\mathrm{JM}}$ 分别为 (FM)、(JM) 的可行解集，
$z^*_{\mathrm{FM}}$、$z^*_{\mathrm{JM}}$ 为相应最优值。在相同数据下
$\mathcal{F}_{\mathrm{FM}}\subseteq\mathcal{F}_{\mathrm{JM}}$，从而 $z^*_{\mathrm{JM}}\le z^*_{\mathrm{FM}}$。
\end{proposition}

\begin{proof}
(FM) 的约束是 (JM) 的约束加上式~\eqref{eq:fixed}，而由式~\eqref{eq:data-fixed}，
式~\eqref{eq:fixed} 给出的维护时刻满足式~\eqref{eq:window}。
所以 (FM) 的每个可行解都是 (JM) 的可行解，两者目标函数相同，故联合最优值不大于固定最优值。
两个最优值都存在，见推论~\ref{cor:horizon}。
\end{proof}

该关系只对\textbf{两者都求得最优}的情形成立。若某一方只得到限时可行解，
两个数值可能逆序，不能据此认为模型违背该命题，也不能把限时可行解称为最优解。
实验中两种方式使用同一组数据、目标和生产约束，唯一区别是是否放开维护时机，
因此二者之差衡量的就是协同决策维护时机带来的收益。

\section{模型性质}

\subsection{左移解与时域上界}

在可行解中，把同一机器上的工件和维护统称为该机器的\textbf{事件}。
由不重叠约束，同一机器上的事件可以按开始时刻排成一个序列。

\begin{definition}[左移解]\label{def:left-shift}
称可行解为\textbf{左移解}，如果每台机器上的每个事件都尽可能早地开始，即
（记 $e^-$ 为同一机器上的前一个事件，首个事件取 $C_{e^-}=0$）：
\begin{itemize}[nosep]
  \item 工件 $j$：$S_j=C_{e^-}$；
  \item 维护 $(i,k)$：联合方式下 $M_{ik}=\max(C_{e^-},L_{ik})$，固定方式下 $M_{ik}=f_{ik}$。
\end{itemize}
\end{definition}

\begin{lemma}[左移占优]\label{lem:left-shift}
对 (JM) 或 (FM) 的任一可行解 $\sigma$，保持每台机器上的事件集合与先后顺序不变，
按定义~\ref{def:left-shift} 依次重新确定开始时刻，得到的 $\sigma'$ 是同一模型的可行解，
且对每个事件，$\sigma'$ 中的开始时刻不晚于 $\sigma$ 中的开始时刻。
特别地，$C'_j\le C_j$ 对所有 $j$ 成立，从而 $\sum_j T'_j\le\sum_j T_j$。
\end{lemma}

\begin{proof}
固定一台机器 $i$，设其事件序列为 $e_1,e_2,\dots$，对位置 $r$ 归纳证明
$\mathrm{start}'(e_r)\le\mathrm{start}(e_r)$，从而 $\mathrm{end}'(e_r)\le\mathrm{end}(e_r)$。
记 $c=\mathrm{end}(e_{r-1})$、$c'=\mathrm{end}'(e_{r-1})$（$r=1$ 时均为 0），由归纳假设 $c'\le c$；
又因 $\sigma$ 中事件不重叠，$c\le\mathrm{start}(e_r)$。
\begin{itemize}[nosep]
  \item $e_r$ 是工件：$\mathrm{start}'(e_r)=c'\le c\le\mathrm{start}(e_r)$。
  \item $e_r$ 是联合方式下的维护 $(i,k)$：$\sigma$ 满足 $M_{ik}\ge L_{ik}$，
        故 $\max(c',L_{ik})\le\max(c,L_{ik})\le M_{ik}$。
  \item $e_r$ 是固定方式下的维护：$\mathrm{start}'(e_r)=f_{ik}=\mathrm{start}(e_r)$。
\end{itemize}
可行性：每个事件都在前一事件结束后开始（固定方式下 $c'\le c\le f_{ik}$），
所以同一机器上不重叠，式~\eqref{eq:jj}、\eqref{eq:jm} 成立；
联合方式下 $M'_{ik}\ge L_{ik}$ 且 $M'_{ik}+q_{ik}\le M_{ik}+q_{ik}\le U_{ik}$，
固定方式下 $M'_{ik}=f_{ik}$，故式~\eqref{eq:window}、\eqref{eq:fixed} 成立；
分配与时长不变，式~\eqref{eq:assign}、\eqref{eq:proc} 成立。
最后，$T_j$ 是 $C_j$ 的非减函数，所以总延期不增加。
\end{proof}

\begin{corollary}[时域上界]\label{cor:horizon}
(JM) 与 (FM) 都有可行解和最优解；并且存在一个最优解，其所有 $S_j$、$C_j$、$M_{ik}$、$T_j$ 都在 $[0,H]$ 内，
其中 $H=\Ustar+P$。因此把这些变量的取值范围限制为 $[0,H]$ 不会损失最优值。
\end{corollary}

\begin{proof}
\emph{可行解存在}：联合方式取 $M_{ik}=L_{ik}$，固定方式取 $M_{ik}=f_{ik}$；
由式~\eqref{eq:data-window}--\eqref{eq:data-fixed} 这些维护合法且互不重叠，且都在 $\Ustar$ 前结束。
把全部工件依次排在机器 1 的最后一次维护之后，即得到可行解。

\emph{上界}：设 $\sigma'$ 为左移解。在机器 $i$ 上，最后一次维护 $(i,K_i)$ 之前的事件都在
$M'_{iK_i}\le U_{iK_i}-q_{iK_i}<\Ustar$ 前结束；这次维护在 $U_{iK_i}\le\Ustar$ 前结束；
其后的工件首尾相接地加工，总时长不超过 $P$，所以最后完工不晚于 $\Ustar+P=H$。
于是左移解的所有时刻都在 $[0,H]$ 内，且 $T_j\le C_j\le H$。

\emph{最优解存在}：左移解由各机器的事件集合与顺序唯一确定，因此只有有限个。
由引理~\ref{lem:left-shift}，任一可行解都不优于某个左移解，所以在有限个左移解中取目标最小者，
就得到全局最优解，并且它满足上述上界。
\end{proof}

\subsection{计算复杂性}

记 (JM) 和 (FM) 的判定形式为：给定实例和整数 $y\ge 0$，是否存在可行解使 $\sum_j T_j\le y$。

\begin{theorem}\label{thm:np}
(JM) 与 (FM) 的判定形式都属于 NP，并且：
\begin{enumerate}[label=(\alph*),nosep]
  \item 即使 $m=1$、$K_1=1$，它们也是 NP 完全的；
  \item 当机器数 $m$ 作为输入的一部分时，即使每台机器只有一个时间窗，它们也是强 NP 完全的。
\end{enumerate}
单窗口模型是多窗口模型的特例，所以结论对两种模型都成立。
\end{theorem}

\begin{proof}
\emph{属于 NP}：由推论~\ref{cor:horizon} 的证明，若存在满足 $\sum_jT_j\le y$ 的可行解，
则存在这样的左移解，其所有时刻不超过 $H$，编码长度是输入规模的多项式；
按式~\eqref{eq:assign}--\eqref{eq:fixed} 逐条验证只需多项式时间。

\emph{(a)} 由单机总延期问题 $1\|\sum T_j$ 归约，该问题的判定形式是（普通意义下）NP 完全的~\cite{du1990}。
给定 $1\|\sum T_j$ 的实例 $(p_j,d_j)_{j\in\cJ}$ 与 $y$，构造：$m=1$，一个时间窗 $[P,P+1)$，
$q=1$，固定方式取 $f=P$，阈值仍为 $y$。下面说明两个问题的最优值相等。
\begin{itemize}[nosep]
  \item 对 $1\|\sum T_j$ 的任一工件顺序，从 0 时刻起无空闲地加工，最后一个工件在 $P$ 完工；
        再在 $[P,P+1)$ 维护，即得到两种方式下的可行解，各工件完工时刻不变。
  \item 反之，对构造实例的任一可行解，保持工件顺序，删去维护并从 0 时刻起无空闲地加工；
        每个工件的完工时刻等于其前面工件（含自身）的加工时间之和，
        不晚于原解中的完工时刻，所以总延期不增加。
\end{itemize}
因此构造实例存在 $\sum T_j\le y$ 的解当且仅当原实例存在，且构造是多项式的。

\emph{(b)} 由 3-划分问题归约，该问题是强 NP 完全的~\cite{garey1979}：
给定正整数 $a_1,\dots,a_{3t}$ 与 $B$，满足 $B/4<a_r<B/2$ 且 $\sum_ra_r=tB$，
问能否把它们分成 $t$ 组，使每组之和都等于 $B$（由取值范围，这样的每组恰有 3 个数）。
构造：$m=t$ 台机器，$n=3t$ 个工件，$p_r=a_r$，$d_r=B$；每台机器一个时间窗 $[tB,\,tB+1)$，
$q=1$，固定方式取 $f=tB$；阈值 $y=0$。
\begin{itemize}[nosep]
  \item 若存在 3-划分，把每组分到一台机器，从 0 时刻起无空闲地加工，每台机器在 $B\le tB$ 时刻完成；
        再在 $[tB,tB+1)$ 维护。所有 $C_r\le B=d_r$，总延期为 0。
  \item 若存在总延期为 0 的可行解，则所有 $C_r\le B$，每台机器上的工件都在 $[0,B)$ 内互不重叠地加工，
        负载不超过 $B$；而总负载为 $tB$，所以每台机器的负载恰为 $B$，即得到一个 3-划分。
\end{itemize}
构造实例中的数值都不超过 $tB+1$，是 3-划分实例一元编码长度的多项式，故归约保持强 NP 完全性。
\end{proof}

定理~\ref{thm:np} 说明，除非 $\mathrm{P}=\mathrm{NP}$，本问题不存在多项式时间精确算法；
机器数作为输入时，连伪多项式时间算法也不存在。
这就是本文在大规模算例上采用调度规则、元启发式和强化学习，
而只在小规模算例上用精确方法作参考的原因。
文献~\cite{lee2000} 对每台机器在 $[0,T]$ 内维护一次、目标为总完工时间的情形证明了 NP 难，
其模型与目标都和本文不同，所以本文给出了上面的独立归约。

\subsection{构造式决策过程的完备性}\label{sec:env}

第3章的调度构造过程建立在下面的逐机器追加方式之上。
状态记录每台机器的当前结束时刻 $\tau_i$（初值 0）、已安排的工件，以及每台机器下一个待做的时间窗 $k_i$。
记机器 $i$ 上“下一个工件最晚完工时刻”为
\[
\lambda_i=
\begin{cases}
U_{ik_i}-q_{ik_i}, & \text{联合方式且机器 $i$ 还有待做窗口},\\
f_{ik_i}, & \text{固定方式且机器 $i$ 还有待做窗口},\\
+\infty, & \text{机器 $i$ 的窗口已全部完成}.
\end{cases}
\]
可选动作有两类：
\begin{itemize}[nosep]
  \item 工件动作 $(j,i)$：未安排的工件 $j$ 在机器 $i$ 的 $\tau_i$ 时刻开工，要求 $\tau_i+p_j\le\lambda_i$；
  \item 维护动作 $i$：机器 $i$ 还有待做窗口时，在 $\max(\tau_i,L_{ik_i})$（联合）或 $f_{ik_i}$（固定）开始维护。
\end{itemize}
每个动作执行后更新 $\tau_i$ 和 $k_i$。所有工件和维护都安排完后，过程结束，共 $n+W$ 步。

\begin{proposition}\label{prop:env}
对上述构造过程：
\begin{enumerate}[label=(\roman*),nosep]
  \item 任何未结束的状态都至少有一个可选动作，每次构造都在 $n+W$ 步后得到一个可行解；
  \item 每个左移解都能由某个动作序列构造出来。
\end{enumerate}
因此，所有构造结果中的最优者就是 (JM)（或 (FM)）的最优解。
\end{proposition}

\begin{proof}
(i) 先证不变式：机器 $i$ 有待做窗口 $k_i$ 时，$\tau_i\le\lambda_i$。
初始 $\tau_i=0\le L_{i1}\le\lambda_i$；追加工件后由动作条件 $\tau_i\le\lambda_i$；
完成维护 $(i,k)$ 后 $\tau_i\le U_{ik}\le L_{i,k+1}\le\lambda_i$。
由该不变式，有待做窗口的机器的维护动作总是合法的：联合方式下
$L_{ik_i}\le\max(\tau_i,L_{ik_i})\le U_{ik_i}-q_{ik_i}$，固定方式下 $\tau_i\le f_{ik_i}$。
若没有任何待做窗口，则所有 $\lambda_i=+\infty$，任一未安排工件都可以追加。
所以未结束的状态总有可选动作；每步恰好安排一个工件或一次维护，故 $n+W$ 步后结束。
每个事件都在本机器当前结束时刻或之后开始，同一机器上不重叠；
维护落在时间窗内（固定方式下等于 $f_{ik}$）；所有工件和时间窗都已覆盖。
因此结果是可行解。

(ii) 设 $\sigma'$ 为左移解。把所有事件按任意一种保持各机器内部先后顺序的方式排成序列
（例如按开始时刻），依次执行对应动作。对执行步数归纳，执行到机器 $i$ 的事件 $e$ 时，
$\tau_i$ 等于 $\sigma'$ 中 $e$ 的前一事件的结束时刻（其他机器的动作不改变 $\tau_i$）。
\begin{itemize}[nosep]
  \item $e$ 是工件 $j$：动作在 $\tau_i=S'_j$ 开工，与 $\sigma'$ 一致。若此时机器 $i$ 还有待做窗口 $k_i$，
        则该维护在 $\sigma'$ 中排在 $j$ 之后，故 $C'_j\le M'_{ik_i}\le\lambda_i$，动作合法。
  \item $e$ 是维护 $(i,k)$：由式~\eqref{eq:data-order}，同一机器的维护按窗口顺序出现，
        所以此时 $k_i=k$；动作的开始时刻 $\max(\tau_i,L_{ik})$（或 $f_{ik}$）正是左移解的定义值。
\end{itemize}
所以动作序列合法，并且逐个复现 $\sigma'$。
结合引理~\ref{lem:left-shift} 与推论~\ref{cor:horizon}，存在左移的最优解，它可被构造出来；
又由 (i)，构造结果都是可行解，所以构造结果中的最优值等于模型最优值。
\end{proof}

\begin{remark}\label{rem:decoder}
第3章的强化学习主方法使用的是另一种构造方式：在元启发式的序列解码器上逐个选择“下一个工件”
或“某台机器的下一次维护”，而工件被放到使其\textbf{完工最早}的机器上。
这一规则限制了机器选择，命题~\ref{prop:env} 的论证不再适用。
没有维护时，同质机器可以互换，经典的列表调度论证保证这种规则不丢失最优解；
但各机器的维护窗口不同，机器不再可以互换，本文\textbf{没有证明}该解码器的序列空间一定包含最优解。
作为经验验证，在 $n=3$、$m=2$ 的 16 个单窗口与 16 个多窗口算例上，
对两种维护方式（共 64 种组合）穷举全部序列，所得最好值都等于 CP-SAT 证得的最优值
（回归测试 \texttt{test\_both\_construction\_spaces\_contain\_a\_proven\_optimum}）。
这只说明在这些小算例上没有发现反例，不能推广为一般结论。
\end{remark}

\section{约束规划模型}\label{sec:cpsat}

小规模算例的精确参考解由 OR-Tools CP-SAT~\cite{ortools} 求得，模型与 (JM)/(FM) 一一对应：
\begin{itemize}[nosep]
  \item 每个工件 $j$ 有整数变量 $S_j,C_j\in[0,H]$；对每台机器 $i$ 建一个可选区间
        $I_{ji}=[S_j,C_j)$，长度 $p_j$，存在当且仅当 $x_{ji}=1$；用 \texttt{ExactlyOne} 表示式~\eqref{eq:assign}。
  \item 每个时间窗 $(i,k)$ 建一个\textbf{必选}区间，开始变量的取值范围为 $[L_{ik},\,U_{ik}-q_{ik}]$，
        即式~\eqref{eq:window}；固定方式另加 $M_{ik}=f_{ik}$。维护区间不能设为可选，否则求解器可以删掉维护。
  \item 每台机器对其所有存在的加工区间与维护区间施加 \texttt{NoOverlap}，即式~\eqref{eq:jj}、\eqref{eq:jm}。
  \item $T_j\in[0,H]$，用 \texttt{AddMaxEquality} 表示 $T_j=\max(0,C_j-d_j)$；目标为 $\min\sum_jT_j$。
\end{itemize}
由推论~\ref{cor:horizon}，取值范围 $[0,H]$ 不损失最优值。

对求解结果做以下核对后才接受：
\begin{enumerate}[nosep]
  \item 从解中取出分配与时刻，重新构造调度，并独立重算总延期 $T$；
  \item 求解器给出的目标值 $z$ 与下界 $b$ 必须为整数，且满足 $0\le b\le T\le z$；
        状态为 \texttt{OPTIMAL} 时还要求 $b=T=z$；
  \item 调度必须通过第~\ref{sec:checker} 节的合法性检查。
\end{enumerate}
只有状态为 \texttt{OPTIMAL} 时才称最优解；状态为 \texttt{FEASIBLE} 时只称限时可行解，
其 $T$ 只是最优值的上界。报告中使用重算的 $T$，而不是求解器的目标值。
在极小算例（$n\le 5$、$m\le 3$）上，CP-SAT 的最优值还与一个独立的穷举程序核对一致。

\section{调度的合法性条件}\label{sec:checker}

所有算法输出的调度都要经过独立的合法性检查器，检查失败就报错，不返回非法调度。
检查器不修复调度，逐条检查：
\begin{enumerate}[nosep]
  \item \textbf{生产部分}（复用 Phase 1 的静态检查器）：每个工件恰好加工一次（式~\eqref{eq:assign}），
        机器编号合法，同一机器上的加工不重叠（式~\eqref{eq:jj}），
        $C_j-S_j=p_j$ 且与实例数据一致（式~\eqref{eq:proc}）；
  \item \textbf{维护部分}：每个时间窗恰好一条维护记录（不能漏、不能重复），
        时长等于 $q_{ik}$，整个区间在 $[L_{ik},U_{ik})$ 内（式~\eqref{eq:window}），
        固定方式下开始时刻等于 $f_{ik}$（式~\eqref{eq:fixed}），
        与本机器上的加工不重叠（式~\eqref{eq:jm}）；
  \item \textbf{目标}：报告的总延期必须等于按 $\sum_j\max(0,C_j-d_j)$ 独立重算的值（式~\eqref{eq:tard}）。
\end{enumerate}
所有区间按左闭右开处理，所有时间必须是整数。
完工时间 $C_{\max}$ 只按生产记录计算，可能早于最后一次维护的结束时刻，这不影响总延期目标。

\section{算例}

下面两个小例子可以手工验证，用来说明联合维护的作用。图中“维”表示维护。

\subsection{单机例子}

两个工件 $p=(2,2)$，$d=(2,4)$；一台机器，维护时长 $q=2$，时间窗 $[0,6)$，固定方式 $f=0$。
\begin{itemize}[nosep]
  \item 固定维护：维护占用 $[0,2)$。顺序 A、B 时 A 在 $[2,4)$、B 在 $[4,6)$，延期 $2+2=4$；
        顺序 B、A 时延期 $0+4=4$。插入空闲只会推迟完工，所以最优总延期为 4。
  \item 联合维护：A 在 $[0,2)$、B 在 $[2,4)$、维护在 $[4,6)$，总延期为 0，达到下界。
\end{itemize}

\begin{figure}[h]
\centering
\begin{tikzpicture}[x=0.75cm,y=0.7cm]
  \node[left,font=\small] at (-0.2,0) {固定：机器};
  \gbar{0}{0}{2}{维}{gray!35}
  \gbar{0}{2}{4}{A}{blue!15}
  \gbar{0}{4}{6}{B}{blue!15}
  \gaxis{0}{6}
  \begin{scope}[xshift=7.2cm]
    \node[left,font=\small] at (-0.2,0) {联合：机器};
    \gbar{0}{0}{2}{A}{blue!15}
    \gbar{0}{2}{4}{B}{blue!15}
    \gbar{0}{4}{6}{维}{gray!35}
    \gaxis{0}{6}
  \end{scope}
\end{tikzpicture}
\caption{单机例子：固定维护总延期 4，联合维护总延期 0}
\end{figure}

\subsection{两机例子}

三个工件 $p=(2,2,2)$，$d=(2,2,4)$；两台机器的维护时长都是 2，时间窗都是 $[0,6)$，固定方式 $f=(0,0)$。
\begin{itemize}[nosep]
  \item 固定维护：两台机器都在 $[0,2)$ 维护，之后三个工件的完工时刻至少为 4、4、6，
        总完工时间至少为 14；交货期之和为 8，而 $\sum_jT_j\ge\sum_j(C_j-d_j)$，所以总延期至少为 6。
        A 在机器 1 的 $[2,4)$、C 在机器 1 的 $[4,6)$、B 在机器 2 的 $[2,4)$ 达到 6，即为最优。
  \item 联合维护：机器 1 加工 A $[0,2)$、C $[2,4)$，然后维护 $[4,6)$；
        机器 2 加工 B $[0,2)$，然后维护 $[2,4)$。三个工件都不延期，总延期为 0。
        此时生产完工时间 $C_{\max}=4$，最后一次维护在 6 结束。
\end{itemize}

\begin{figure}[h]
\centering
\begin{tikzpicture}[x=0.75cm,y=0.7cm]
  \node[left,font=\small] at (-0.2,0) {固定：机器1};
  \node[left,font=\small] at (-0.2,-1) {机器2};
  \gbar{0}{0}{2}{维}{gray!35}
  \gbar{0}{2}{4}{A}{blue!15}
  \gbar{0}{4}{6}{C}{blue!15}
  \gbar{1}{0}{2}{维}{gray!35}
  \gbar{1}{2}{4}{B}{blue!15}
  \gaxis{1}{6}
  \begin{scope}[xshift=7.2cm]
    \node[left,font=\small] at (-0.2,0) {联合：机器1};
    \node[left,font=\small] at (-0.2,-1) {机器2};
    \gbar{0}{0}{2}{A}{blue!15}
    \gbar{0}{2}{4}{C}{blue!15}
    \gbar{0}{4}{6}{维}{gray!35}
    \gbar{1}{0}{2}{B}{blue!15}
    \gbar{1}{2}{4}{维}{gray!35}
    \gaxis{1}{6}
  \end{scope}
\end{tikzpicture}
\caption{两机例子：固定维护总延期 6，联合维护总延期 0}
\end{figure}

两个例子在两种维护方式下的最优值（4 与 0、6 与 0）均已由 CP-SAT（状态 \texttt{OPTIMAL}）
和独立穷举程序复核一致。检查器另有回归测试，确认漏维护、重复维护、维护越出时间窗、
维护与加工重叠、维护时长不符、固定时刻不符、生产部分非法和总延期报错等非法调度都会被拒绝。

\section{本章小结}

本章建立了带周期性预防维护时间窗的同质平行机总延期调度模型，
给出了固定维护与联合维护两种方式，以及单窗口这一特例。主要结论：
\begin{enumerate}[nosep]
  \item 固定维护的可行域是联合维护的子集，联合最优总延期不大于固定最优总延期（命题~\ref{prop:fixed-joint}）；
  \item 总存在左移的最优解，所有时刻都不超过 $H=\Ustar+P$（引理~\ref{lem:left-shift}、推论~\ref{cor:horizon}）；
  \item 问题即使只有一台机器也是 NP 难的，机器数作为输入时是强 NP 难的（定理~\ref{thm:np}）；
  \item 逐机器追加的构造过程不会陷入死局，且其构造空间包含最优解（命题~\ref{prop:env}）；
        强化学习所用的“完工最早机器”解码器是否也包含最优解尚未证明，只在小算例上穷举验证过。
\end{enumerate}
这些结论依赖本章的假设：同质机器、单工序、0 时刻到达、正则目标、不同机器维护互不制约。
若引入释放时间、工艺路线、换型、共享维修人员或非正则目标，需要重新论证。

\begin{thebibliography}{9}
\bibitem{lee2000}
C.-Y. Lee, Z.-L. Chen. Scheduling jobs and maintenance activities on parallel machines.
\emph{Naval Research Logistics}, 2000, 47(2): 145--165.

\bibitem{du1990}
J. Du, J. Y.-T. Leung. Minimizing total tardiness on one machine is NP-hard.
\emph{Mathematics of Operations Research}, 1990, 15(3): 483--495.

\bibitem{garey1979}
M. R. Garey, D. S. Johnson. \emph{Computers and Intractability: A Guide to the Theory of NP-Completeness}.
San Francisco: W. H. Freeman, 1979.

\bibitem{ortools}
L. Perron, F. Didier. CP-SAT. Google OR-Tools, \url{https://developers.google.com/optimization}.
\end{thebibliography}

\end{document}
